\chapter*{Acknowledgements}
\addcontentsline{toc}{chapter}{Acknowledgements}
\markboth{Acknowledgements}{Acknowledgements}

I would like to thank the lecturers of the Faculty of Computer Science and
Engineering, Ho Chi Minh City University of Technology (HCMUT), for the
foundational and specialised teaching that made this work possible.

I am especially grateful to my instructor, \InstructorName, for the guidance
that shaped the scope of this work. Two pieces of that guidance did most of the
work. The first was the insistence that a research question be stated together
with the outcome that would \emph{refute} it; applying that test removed several
items from an earlier list that turned out to be names for groups of experiments
rather than questions. The second was the instruction to measure the resource
the system actually spends rather than the one the literature happens to report,
which is why Chapter~\ref{ch:experiments} reports CPU seconds and test
executions side by side instead of either alone.

I also thank my colleagues for the discussions that sharpened the central
distinction this work is built on: that a repair loop learns two different
things on every iteration --- which test is worth running, and which repair idea
is already dead --- and that both are usually discarded.

This document reports work in progress. The implementation and its unit tests
are complete; the $100$-task pilot is running at the time of writing, and
Chapter~\ref{ch:results} marks every number that depends on it. Any remaining
errors are my own.
